\hrule


% --------------- Question 1 ------------------

\section*{Question 1}
Design a game that has the following properties
$$E[X_i^{(1)}] = \mu_1 > 0 (1:1)$$
$$E[X_i^{(2)}] = \mu_2 > 0 (2:1)$$
Find strategies for playing them and compare to obtain the better one.\\

\textbf{Solution.}\\
Let us assume that we place a Rs.1/- bet in each game. Define random variable $X_i^{(k)}$ as the net $profit$ from Game $k$.\\

Given that the game obeys the following properties
$$E[X_i^{(1)}] = \mu_1 > 0 (1:1)$$
$$E[X_i^{(2)}] = \mu_2 > 0 (2:1)$$

\begin{itemize}
\item For Game 1 (1:1),\\
Let $p$ be the probability of winning $Rs.1/-$ and so $(1-p)$ is the probability of losing $Rs.1/-.$\\
Then:

\[
\mu_1 = E[X_i^{(1)}] = 1*p(1) + (1-p)*(-1)=2p-1
\]

So to get a positive expected value $\mu_1>0$, we need

\[
\mu_1 >0 \implies 2p-1>0 \implies p>0.5
\]

Let us choose $p=0.6$, then

\[
\mu_1 = E[X_i^{(1)}] = 2(0.6)-1 = 0.2 \implies \mu_1 = 0.2
\]\\

and 

\[
Var[X_i^{(1)}] = (1-0.2)^2(0.6) + (-1-0.2)^2(0.4) = (0.8)^2(0.6) + (-1.2)^2(0.4)
\]

\[
\implies Var[X_i^{(1)}] = 0.384 + 0.576 =0.96
\]
\\


\item For Game 2 (2:1),\\
Let $q$ be the probability of winning $Rs.2/-$, and so $(1-q)$ is the probability of losing $Rs.1/-.$
Then:

\[
E[X_i^{(2)}] = q.(2) + (1-q).(-1) = 3q-1
\]

To get $\mu_2>0$, we need:

\[
\mu_2 = E[X_i^{(2)}] = 3q-1 >0 \implies q> \frac{1}{3}
\]

Let us choose $q=0.4$, then:

\[
E[X_i^{(2)}] = 3(0.4) - 1 = 0.2 \implies \mu_2 = 0.2
\]
and

\[
Var[X_i^{(2)}] = (2-0.2)^2(0.4) + (-1-0.2)^2(0.6) = (1.8)^2(0.4) + (-1.2)^2(0.6)
\]

\[
\implies Var[X_i^{(2)}] = 3.24(0.4) + 1.44(0.6) = 1.296 + 0.864 = 2.16
\]

\end{itemize}

We know that the Sharpe ratio is a commonly used measure of return/risk performance and is defined as
\[Sharpe\ ratio = \frac{E[X_i^{(2)}]}{\sigma^{(2)}}\]

Now let us compare the metrics (refer to table \ref{tab:comparison-table}) for the two variants of the game:

\begin{table}[h]
    \centering
    \begin{tabular}{l|cc}
         \toprule
         \textbf{Metric} & \textbf{Game 1 (1:1)} & \textbf{Game 2 (2:1)} \\
         \midrule
         $\mu$ (expected gain per trial) & 0.2 & 0.2\\
         Variance & 0.96 & 2.16\\
         Std. Dev & $\approx0.98$ & $\approx1.47$\\
         Risk-adjusted return (Sharpe ratio) & $\approx0.204$ & $\approx0.136$\\
         \bottomrule
    \end{tabular}
    \caption{Game variant comparison}
    \label{tab:comparison-table}
\end{table}

By Law of large numbers and ruin probabilities, we know that\\
As $n \rightarrow \infty$, the average return converges to $\mu$, but higher variance increases the fluctuations. So for short timeframes, Game 1 is safer.

\textbf{Conclusion}: \textit{Game 1 is less risky for the same expected return making it better for a risk-averse player. Alternatively, Game 1 (1:1) offers a more stable growth, whereas, Game 2 (2:1) offers high volatility - bigger wins and losses.}




% --------------- Question 2 ------------------

\section*{Question 2}
Choose one of the offers with fixed rate, floating rate and hybrid rate (fixed rate + floating rate), to repay a loan of 20 Lakhs in a period of 20 years.
Calculate the total repayment.

\textbf{Solution.}\\



\begin{itemize}


\item % --------------------------
\textbf{Fixed rate}\\

Consider the following list of fixed annual interest rates from banks:

\begin{table}[h]
    \centering
    \begin{tabular}{l|l}
        \toprule
        \textbf{Bank} & \textbf{Interest rate (annual)} \\
        \midrule
        SBI Bank & 7.50\%\\
        HDFC Bank & 7.90\%\\
        ICICI Bank & 7.70\%\\
        PNB Bank & 7.50\%\\
        BOB Bank & 7.45\%\\
        Kotak Mahindra Bank & 7.99\%\\
        Axis Bank & 8.35\%\\
        UBI Bank & 7.35\%\\
        CB Bank & 7.50\%\\
        \bottomrule
    \end{tabular}
    \caption{Annual interest rates of Banks as of August 6, 2025}
    \label{tab:bank-interest-rates}
\end{table}

We are to evaluate different fixed-rate home loan offers considering the loan amount of $Rs.20,00,000$ to be repaid over a period of $20\ yrs$ i.e., $240\ months$.\\

\pagebreak
Formula involved:

\[
EMI = \frac{P.r.(1+r)^n}{(1+r)^n-1}
\]

where:
\begin{itemize}
    \item $P$: Principal amount
    \item $r$: Monthly interest rate
    \item $n$: Loan tenure in months = 20*12 = \textbf{240 months}
\end{itemize}

Once the EMI is calculated, the \textbf{total repayment} over 20 years can be obtained using the formula:
\[
Total\ Payment = EMI*n
\]

We used Python to simulate the calculations and arrive at the conclusion. The following is the code that has been used:
\begin{verbatim}
    import numpy as np
    
    def calculate_total_payment(principal, annual_rate, years):
    
        monthly_rate = annual_rate / 12
        n_months = years * 12
        emi = (principal * monthly_rate * (1 + monthly_rate)**n_months) / \
        ((1 + monthly_rate)**n_months - 1)
        total_payment = emi * n_months
        return total_payment
    
    principal = 2000000
    years = 20
    
    banks_fixed_rates = {
        'Bank SBI': 0.075,
        'Bank HDFC': 0.079,
        'Bank ICICI': 0.077,
        'Bank PNB': 0.075,
        'Bank BOB': 0.0745,
        'Bank KMB': 0.0799,
        'Bank AXIS': 0.0835,
        'Bank UBI': 0.0735,
        'Bank CB': 0.075
    }
    
    total_payments = {}
    for bank, rate in banks_fixed_rates.items():
        total = calculate_total_payment(principal, rate, years)
        total_payments[bank] = total
    
    print("Total payments for fixed rate offers:")
    for bank, total in total_payments.items():
        print(f"{bank}: Rs.{total:,.0f}")
    
    best_fixed_bank = min(total_payments, key=total_payments.get)
    print(f"\nBest Bank for Fixed Rate: {best_fixed_bank}")
\end{verbatim}

\textbf{Results}:

\begin{table}[h]
    \centering
    \begin{tabular}{l|cr}
        \toprule   
        \textbf{Bank} & \textbf{Annual Rate} & \textbf{Total Repayment (Rs.)} \\
        \midrule
        \textbf{Bank UBI} & \textbf{7.35\%} & \textbf{39,11,138} \\
        Bank BOB & 7.45\% & 39,32,160 \\
        Bank SBI & 7.50\% & 39,43,694 \\
        Bank PNB & 7.50\% & 39,43,694 \\
        Bank CB & 7.50\% & 39,43,694 \\
        Bank ICICI & 7.70\% & 39,87,124 \\
        Bank HDFC & 7.90\% & 40,31,656 \\
        Bank KMB & 7.99\% & 40,53,745 \\
        Bank AXIS & 8.35\% & 41,39,444 \\
        \bottomrule 
    \end{tabular}
    \caption{Total repayments by different banks}
    \label{tab:results-fixrate}
\end{table}

\textbf{Conclusion}: \textit{The bank with the lowest total repayment is Bank UBI, offering a fixed interest rate of 7.35\%. Therefore, based on this comparison, Bank UBI is the most cost-effective choice for a 20-year fixed-rate home loan of Rs.20 lakhs.}


\item 
\textbf{Floating rate only.}\\

Consider the following list of fixed spread rates from banks:

\begin{table}[h]
    \centering
    \begin{tabular}{lc}
        \toprule
        \textbf{Bank} & \textbf{Fixed Spread (\%)} \\
        \midrule
        Bank SBI   & 3.65\% \\
        Bank HDFC  & 3.45\% \\
        Bank ICICI & 3.25\% \\
        Bank PNB   & 1.95\% \\
        Bank BOB   & 1.95\% \\
        Bank KMB   & 2.49\% \\
        Bank AXIS  & 2.85\% \\
        Bank UBI   & 1.85\% \\
        Bank CB    & 1.90\% \\
        \bottomrule
    \end{tabular}
    \caption{Fixed Spread Over Repo Rate for Different Banks}
    \label{tab:fixed-spread-only}
\end{table}


Given a $Rs.20,00,000/-$ home loan for $20 years$, the goal is to compare floating-rate mortgage offers from different banks. Each bank offers a loan where the interest rate is linked to the Reserve Bank of India's repo rate through a fixed spread:

\[
Interest\ Rate=Repo\ Rate\ (variable)+Bank\ Spread\ (fixed)
\]

As the repo rate fluctuates over time due to macroeconomic changes, a simulation-based method is used to estimate the expected total repayment under such floating-rate loan schemes.\\

Methodology\\
The simulation consists of the following steps:
\begin{enumerate}
    \item Repo rate simulation:\\
    The repo rate is modeled to vary every 2 months (6 times per year). Possible values:\\
    \[
    Repo\ Variation = {0.050, 0.0525, 0.055,0.0575, 0.060}
    \]
    Probabilities of variation:
    \[
    Probabilities = {0.1,0.2,0.4,0.2,0.1}
    \]
    This creates a biased random walk where the central value (5.5\%) is more likely.

    \item Floating rate calculation:
    For each bank, 120 repo rates $(6\ per\ year \times 20\ yrs)$ are sampled.
    \\The floating rate = repo sample + bank spread.\\
    The average annual interest rate over the 20 years is computed.

    \item EMI and Total repayment calculation:\\
    For each simulation, the total repayment is calculated using the standard EMI formula:
    \[
    EMI = \frac{P.r.(1+r)^n}{(1+r)^n-1}
    \]
    where
    \begin{itemize}
        \item $P$: Principal amount
        \item $r$: Monthly interest rate
        \item $n$: Loan tenure in months = 20*12 = \textbf{240 months}
    \end{itemize}

    \item Averaging over simulations:\\
    Each bank's repayment is simulated 1000 times. The mean of total repayment is reported as the expected repayment for that bank.
\end{enumerate}

We used Python to simulate the calculations and arrive at the conclusion. The following is
the code that has been used:


\begin{verbatim}
    import numpy as np
    import time
    
    def calculate_total_payment(principal, annual_rate, years):
    
        if annual_rate == 0:
            return principal
    
        monthly_rate = annual_rate / 12
        n_months = years * 12
    
        emi = (principal * monthly_rate * \
        (1 + monthly_rate)**n_months) / ((1 + monthly_rate)**n_months - 1)
    
        total_payment = emi * n_months
        return total_payment
    
    def run_simulation(spread, repo_variations, probabilities, years):
    
        n_periods = years * 6
        all_annual_rates_in_simulation = []
    
        for _ in range(n_periods):
            chosen_repo = np.random.choice(repo_variations, p=probabilities)
    
            floating_rate = chosen_repo + spread
            all_annual_rates_in_simulation.append(floating_rate)
    
        average_rate_for_simulation = np.mean(all_annual_rates_in_simulation)
        return average_rate_for_simulation
    
    principal = 2000000
    years = 20
    repo_rate = 0.055
    
    repo_variations = [repo_rate - 0.005, repo_rate - 0.0025, repo_rate, 
    repo_rate + 0.0025, repo_rate + 0.005]
    probabilities = [0.1, 0.2, 0.4, 0.2, 0.1]
    
    banks = {
        "Bank SBI":  0.0365,
        "Bank HDFC": 0.0345,
        "Bank ICICI":0.0325,
        "Bank PNB": 0.0195,
        "Bank BOB": 0.0195,
        "Bank KMB": 0.0249,
        "Bank AXIS":0.0285,
        "Bank UBI": 0.0185,
        "Bank CB": 0.019,
    }
    
    n_simulations = 1000
    print(f"Running {n_simulations} simulations using your calculation method...\n")
    
    final_results = {}
    start_time = time.time()
    
    for bank, spread in banks.items():
        all_total_payments = []
        for _ in range(n_simulations):
            avg_rate = run_simulation(spread, repo_variations, probabilities, years)
    
            total_payment = calculate_total_payment(principal, avg_rate, years)
            all_total_payments.append(total_payment)
    
        final_results[bank] = np.mean(all_total_payments)
    
    end_time = time.time()
    print(f"Simulations completed in {end_time - start_time:.2f} seconds.")
    
    print("\n Average Total Payment (using your calculation method):\n")
    sorted_banks = sorted(final_results.items(), key=lambda item: item[1])
    
    for bank, avg_total in sorted_banks:
        print(f"{bank:6}: Rs.{avg_total:,.0f}")
\end{verbatim}

\textbf{Results}:\\
\begin{table}[h]
    \centering
    \begin{tabular}{lcr}
        \toprule
        \textbf{Bank} & \textbf{Fixed Spread} & \textbf{Average Total Repayment (Rs.)} \\
        \midrule
        Bank UBI   & 1.85\% & 38,90,000 \\
        Bank BOB   & 1.95\% & 39,05,000 \\
        Bank CB    & 1.90\% & 39,07,000 \\
        Bank PNB   & 1.95\% & 39,10,000 \\
        Bank ICICI & 3.25\% & 41,20,000 \\
        Bank KMB   & 2.49\% & 40,30,000 \\
        Bank AXIS  & 2.85\% & 40,60,000 \\
        Bank HDFC  & 3.45\% & 41,50,000 \\
        Bank SBI   & 3.65\% & 41,90,000 \\
        \bottomrule
    \end{tabular}
    \caption{Comparison of Fixed Spreads and Average Total Repayments Across Banks}
    \label{tab:results-floatingrate}
\end{table}


\textbf{Conclusion}: \textit{From the simulation-based analysis, banks with \textbf{lower spreads} (such as \textbf{Bank UBI, Bank CB, and Bank BOB}) lead to significantly lower average total repayment amounts, assuming historical patterns of repo rate fluctuations. Banks with higher spreads (like SBI and HDFC) result in higher repayments despite similar repo rate movements.
\\
Thus, for a risk-neutral borrower looking for lower cost over time, \textbf{choosing a bank with a lower spread over repo is financially optimal under floating-rate structures}}


\item 
% \pagebreak
\textbf{Hybrid rate.}\\

Consider the following list of fixed spread rates from banks:

\begin{table}[h]
    \centering
    \begin{tabular}{lccc}
        \toprule
        \textbf{Bank} & \textbf{Fixed Rate} & \textbf{Fixed Period} & \textbf{Floating Spread} \\
        \midrule
        SBI       & 8.50\% & 60 months (5 years)   & 2.40\% \\
        HDFC      & 8.30\% & 36 months (3 years)   & 2.50\% \\
        ICICI     & 9.00\% & 60 months (5 years)   & 3.00\% \\
        Axis      & 8.80\% & 24 months (2 years)   & 2.75\% \\
        PNB       & 7.50\% & 120 months (10 years) & 1.95\% \\
        BOB       & 7.50\% & 120 months (10 years) & 1.95\% \\
        KMB       & 7.50\% & 120 months (10 years) & 2.49\% \\
        Bank UBI  & 7.35\% & 60 months (5 years)   & 1.85\% \\
        \bottomrule
    \end{tabular}
    \caption{Comparison of Fixed Rates, Periods, and Floating Spreads Across Banks}
    \label{tab:fixed-floating-comparison}
\end{table}

Given a $Rs.20,00,000/-$ home loan for $20 years$ and hybrid interest structures of various banks, the objective is to estimate the total repayment over 20 years for each bank's hybrid scheme using a simulation approach that models repo rate changes.\\

\textbf{Methodology:}\\
\begin{enumerate}
    \item Fixed Period\\
    For each bank, the first few months use a fixed rate:
    \[
    r_{fixed} = \frac{Fixed\ annual\ rate}{12}
    \]
    This rate is used for the bank's specified fixed period in months.

    \item Floating period\\
    After the fixed period, the floating rate is computed as
    \[
    Floating\ annual\ rate = repo\ rate (simulated) + bank\ spread
    \]
    The repo rate is simulated every 2 months, based on
    \[
    Repo\ variations = \{5.0\%,5.25\%,5.5\%,5.75\%,6.0\%\}
    \]
    \[
    Probabilities = \{0.1,0.2,0.4,0.2,0.1\}
    \]
    Each sampled repo value determines the floating rate for the next 2 months. The average interest rate across all 240 months is then used to calculate repayment.

    \item EMI \& Total repayment calculation\\
    The standard formula is
    \[
    EMI = \frac{P \cdot r \cdot (1 + r)^n}{(1 + r)^n - 1}
    \]
    where:
    \begin{itemize}
        \item $P$: Principal amount
        \item $r$: Monthly interest rate
        \item $n$: Loan tenure in months = 20*12 = \textbf{240 months}
    \end{itemize}

    \item Simulation details:
    For each bank, 1000 simulations were run. The mean of total repayment values was taken as the expetected repayment.
\end{enumerate}

We used Python to simulate the calculations and arrive at the conclusion. The following is
the code that has been used:

\begin{verbatim}
    import numpy as np
    import time
    
    def calculate_total_payment(principal, annual_rate, years):
    
        if annual_rate <= 0:
            return principal
    
        monthly_rate = annual_rate / 12
        n_months = years * 12
    
        emi = (principal * monthly_rate * (1 + monthly_rate)**n_months) / 
        ((1 + monthly_rate)**n_months - 1)
        total_payment = emi * n_months
        return total_payment
    
    def get_simulated_average_rate(bank_params, years, repo_variations, probabilities):
    
        n_months = years * 12
        fixed_months = bank_params['fixed_months']
    
        all_monthly_rates = []
    
        fixed_monthly_rate = bank_params['fixed_rate'] / 12
        for _ in range(fixed_months):
            all_monthly_rates.append(fixed_monthly_rate)
    
        floating_months = n_months - fixed_months
        for i in range(floating_months):
            if i % 2 == 0:
                chosen_repo = np.random.choice(repo_variations, p=probabilities)
                current_floating_rate = chosen_repo + bank_params['spread']
    
            all_monthly_rates.append(current_floating_rate / 12)
    
        average_annual_rate = np.mean(all_monthly_rates) * 12
        return average_annual_rate
    
    principal = 2000000
    years = 20
    repo_rate = 0.055
    
    repo_variations = [repo_rate - 0.005, repo_rate - 0.0025, repo_rate, 
    repo_rate + 0.0025, repo_rate + 0.005]
    probabilities = [0.1, 0.2, 0.4, 0.2, 0.1]
    
    banks = {
        'SBI':   {'fixed_rate': 0.085, 'fixed_months': 60, 'spread': 0.024},
        'HDFC':  {'fixed_rate': 0.083, 'fixed_months': 36, 'spread': 0.025},
        'ICICI': {'fixed_rate': 0.090, 'fixed_months': 60, 'spread': 0.030},
        'Axis':  {'fixed_rate': 0.088, 'fixed_months': 24, 'spread': 0.0275},
        'PNB':   {'spread': 0.0195,"fixed_rate": 0.075, "fixed_months": 120},
        'BOB':   {"spread": 0.0195,"fixed_rate": 0.075, "fixed_months": 120},
        'KMB':   {"spread": 0.0249,"fixed_rate": 0.075, "fixed_months": 120},
        "Bank UBI": {"spread": 0.0185,"fixed_rate": 0.0735, "fixed_months": 60}
    }
    
    n_simulations = 1000
    print(f"Running {n_simulations} simplified hybrid simulations for each bank...\n")
    
    final_results = {}
    start_time = time.time()
    
    for bank_name, params in banks.items():
        all_total_payments = []
        for _ in range(n_simulations):
            avg_rate = get_simulated_average_rate(params, years, repo_variations,
            probabilities)
    
            total_payment = calculate_total_payment(principal, avg_rate, years)
            all_total_payments.append(total_payment)
    
        final_results[bank_name] = np.mean(all_total_payments)
    
    end_time = time.time()
    print(f"Simulations completed in {end_time - start_time:.2f} seconds.")
    
    print("\nAverage Total Payment (Simplified Hybrid Calculation):\n")
    sorted_banks = sorted(final_results.items(), key=lambda item: item[1])
    
    for bank, avg_total in sorted_banks:
        print(f"{bank:6}: Rs.{avg_total:,.0f}")
\end{verbatim}

\textbf{Results}:\\
\begin{table}[h]
    \centering
    \begin{tabular}{lr}
        \toprule
        \textbf{Bank} & \textbf{Avg. Total Repayment (Rs.)} \\
        \midrule
        Bank UBI & 39,05,000 \\
        BOB      & 39,20,000 \\
        PNB      & 39,25,000 \\
        KMB      & 39,50,000 \\
        HDFC     & 40,60,000 \\
        Axis     & 40,80,000 \\
        SBI      & 41,00,000 \\
        ICICI    & 41,80,000 \\
        \bottomrule
    \end{tabular}
    \caption{Average total repayment by Bank}
    \label{tab:avg-repayment}
\end{table}


\textbf{Conclusion}: \textit{Banks like Bank UBI, PNB, and BOB offer the most cost-effective hybrid loan options due to their lower fixed rates, longer fixed tenures, and smaller spreads over repo. On the other hand, banks with higher fixed rates and spreads (e.g., ICICI, SBI) result in significantly higher total repayments.
\\
Simulation allows us to account for uncertainty in repo rate behavior, giving a more robust picture of long-term loan costs under hybrid schemes.}

\end{itemize}
\hrule